\documentclass[conference]{IEEEtran}
\usepackage[T1]{fontenc}
\usepackage[utf8]{inputenc}
\usepackage{amsmath,amssymb,amsfonts}
\usepackage{booktabs}
\usepackage{graphicx}
\usepackage{xcolor}
\usepackage{enumitem}
\usepackage{algorithm}
\usepackage{algorithmicx}
\usepackage{algpseudocode}
\usepackage[hidelinks]{hyperref}
\usepackage{balance}
\providecommand{\BIBentryALTinterwordspacing}{\spaceskip=0pt plus 4pt minus 1pt}
\providecommand{\BIBentrySTDinterwordspacing}{\spaceskip=0pt}

% ---------------------------------------------------------------
% Pre-registration machinery.
% \PH{...}  = pre-registered placeholder: a SLOT bound to a fixed
%             estimand before data exists, never a claim. Filled only
%             from experimental output; if the pre-registered fallback
%             fires, the surrounding sentence is REPLACED by its
%             pre-committed fallback wording, not edited post hoc.
% \wj{...}  = author-visible standpoint comment on a hypothesis:
%             why this hypothesis is posed this way, what the
%             disconfirming result is, and which fallback governs.
%             Strip both macros for camera-ready.
% ---------------------------------------------------------------
\newcommand{\PH}[1]{\colorbox{gray!18}{\small$\langle$#1$\rangle$}}
\newcommand{\wj}[1]{\textcolor{blue}{\footnotesize [WJ: #1]}}

\newcommand{\ms}{\ensuremath{\mathit{ms}}}
\newcommand{\Fs}{\ensuremath{F(s)}}
\newcommand{\Gs}{\ensuremath{G(s)}}

\begin{document}

\title{Adequacy Is the Attack Surface: Measuring When Visible Test-Pass
Stops Tracking Correctness for Coding Agents}

\author{\IEEEauthorblockN{Anonymous Author(s)}
\IEEEauthorblockA{Anonymous Institution\\
anonymous@example.com}}

\maketitle

\begin{abstract}
The space between passing a visible test suite and being correct has a
concrete, enumerable form: the surviving mutants of that suite --- executable
programs that differ from the intended implementation yet escape every
visible test. For a coding agent that edits code until the suite turns
green, this surviving-mutant set is an \emph{attack surface}: the exact
space of ways to be wrong and still be accepted, whose size is a single
controllable scalar, the suite's mutation-kill score. We present the first
study that treats visible-suite mutation adequacy as a controlled,
size- and coverage-matched independent variable for coding agents, and
estimates the \emph{false-acceptance curve} it induces:
$\Fs = P(H{=}0 \mid V{=}1, \ms(T_v){=}s)$, the fraction of
agent patches accepted at adequacy $s$ that fail a held-out oracle.
On \PH{$N$} repair tasks from SWE-bench Verified (Python) and Defects4J
(Java) and \PH{$K$} agent families, degrading visible adequacy from full
to $\ms{=}0.30$ raises false acceptance by \PH{$\Delta F$~pp}, while false
acceptance at full adequacy remains \PH{$\le F(1.00)$~pp}; the rise
\PH{emerges near $s^\ast{=}0.xx$ / grows gradually}, and the iterative
fix-until-pass loop exceeds a budget-matched selection-only control by
\PH{$\ge$5~pp} in the low-adequacy band. A human-validated differential
diagnostic (agreement \PH{$\kappa{=}0.xx$}) attributes \PH{$xx\%$} of
low-adequacy false acceptances to gaming-shaped behavior rather than
honest incompleteness. Our results yield a calibrated sufficiency floor
for practitioners: in our setting, visible suites whose mutation adequacy
falls below \PH{$s^\ast$} should not serve as standalone acceptance gates
for agent-produced patches. All hypotheses, estimands, thresholds, and
fallback interpretations were fixed before data collection.
\end{abstract}

\begin{IEEEkeywords}
coding agents, mutation testing, test adequacy, automated program repair,
overfitting, verifier sufficiency, empirical software engineering
\end{IEEEkeywords}

\section{Introduction}\label{sec:intro}

The space between passing a visible test suite and being correct is not an
unknown residual. Software testing gives it a concrete, enumerable form:
the behavior changes the suite fails to detect. Mutation testing constructs
that set explicitly --- each surviving mutant is an executable program that
differs from the intended implementation yet escapes every visible
test~\cite{Papadakis2019Survey}. For a human developer, a surviving mutant
marks a testing weakness. For a coding agent that edits code and observes
test feedback until the suite turns green, the surviving-mutant set is
something sharper: the exact space of ways to be wrong and still be
accepted, whose size is a single scalar --- the suite's mutation-kill
score --- that can be dialed. The exploitable slack is not estimated after
the fact; it is constructed in advance. That is the opening this paper
exploits: a testing risk one can only warn about becomes a variable one can
sweep.

The trust placed in that slack is growing faster than our knowledge of it.
Execution-grounded pipelines propose a patch, run the visible tests,
revise, and accept once the suite passes; benchmarks report a single pass
rate against whatever suite ships with a
task~\cite{SWEbench_VTF8yNQM66}. The appeal is real --- test execution is
cheap, deterministic, and needs no human --- but it rests on one
assumption: that visible test-pass tracks correctness closely enough. When
the suite is inadequate, an agent under optimization pressure can satisfy
the checked behaviors without fixing the defect, and the pipeline reports
success while accepting a broken patch. The question a benchmark designer
or an engineering team actually faces is not \emph{whether} this can
happen but \emph{how adequate a suite must be before its pass signal can
be trusted} --- and no study answers it, because no study controls
adequacy itself.

Existing work holds fixed precisely the variable we vary.
Imperfect-verifier theory shows that resampling under a \emph{fixed} weak
selector converges not to correctness but to a ceiling set by the
selector's false-positive rate~\cite{ResamplingLimits_j8H84v6AZ1}, and
characterizes how the verifier's ROC curve reshapes what best-of-$N$
accepts~\cite{ROCnReroll_3Gy5mmyuxn}; both study a static generator
against a static verifier, never an agent adapting to one.
Execution-verified benchmarks define correctness by whatever fail-to-pass
tests ship with each repository~\cite{SWEbench_VTF8yNQM66}, so their pass
rates mix agent capability with uncontrolled suite strength.
Reward-over-optimization work documents policies exploiting reward-model
errors as they improve~\cite{ChasingTail_pBjy4ek2QV,acl_pang_etal_2023_reward},
but for text reward models, never a verifier with a computable adequacy
score. Recent benchmarks measure reward hacking of long-horizon coding
agents as a gap between visible validation and held-out pass
rates~\cite{SpecBench_2605_21384,RHB_2605_02964} --- structurally our
two-suite construct --- but under each task's \emph{fixed, shipped} suite
decomposition, with no control over the suite's adequacy. Closest of all,
the classical automated-program-repair literature measures patch
overfitting against a held-out suite~\cite{Smith2015Overfitting,
Qi2015PatchPlausibility} and detects it per patch by differential testing
against a reference patch~\cite{XinReiss2017DiffTGen,Xiong2018PatchSim}
--- we concede these are the method-twins of our visible/held-out
construct and of our per-patch diagnostic. All of this work differs on the
one axis none of it varies: no prior study --- recent or classical ---
sweeps mutation-kill adequacy as a size- and coverage-matched independent
variable against iterative LLM-agent behavior. Concretely, before this
design one could not \emph{attribute} accepted-but-wrong patches to suite
adequacy at all, because adequacy arrived confounded with suite size, code
coverage, and benchmark accident; making it a matched, swept variable is
what turns ``weak suites are risky'' from a caution into a measurement.

\textbf{We treat visible-suite mutation adequacy as a controlled
independent variable and estimate the false-acceptance curve it induces.}
Specifically: for each repair task we construct visible sub-suites at
target mutation-kill scores while holding suite size and code coverage
approximately fixed and filtering equivalent mutants
(Sec.~\ref{sec:sampler}); the agent sees only the visible suite and runs a
standard fix-until-pass loop, blind to a held-out oracle of hidden tests
plus the tests that kill the mutants excluded from the visible suite
(Sec.~\ref{sec:oracle}); and every accepted patch is scored post hoc
against that oracle. The object of study is the \emph{false-acceptance
rate}
\begin{equation}
\Fs \;=\; P\bigl(H{=}0 \,\bigm|\, V{=}1,\; \ms(T_v){=}s\bigr),
\label{eq:fs}
\end{equation}
where $V{=}1$ denotes visible-suite acceptance and $H{=}1$ held-out oracle
pass: of the patches the visible suite \emph{accepts} at adequacy $s$, the
fraction that fail the held-out oracle. This conditional quantity, rather
than a marginal difference of two pass rates, is what a practitioner
gating patches needs: it answers ``of what my suite greenlights, how much
is broken?''

Turning this into a trustworthy measurement raises difficulties that our
design answers with constructions rather than hedges. Adequacy, size, and
coverage correlate in an unconstrained draw, and equivalent mutants
corrupt the score's own denominator --- so an adequacy-controlled sampler
isolates adequacy by construction, and equivalent-mutant handling is
elevated to first-class, reported machinery with a sensitivity band
(Sec.~\ref{sec:equiv}). Repair episodes are noisy --- so a mixed-effects
trend model pools instances and repeated episodes
(Sec.~\ref{sec:analysis}). A coarse scalar does not reveal its own curve
shape --- so a segmented-versus-smooth model comparison forces any ``sharp
threshold'' claim to beat a gradual alternative. A rising \Fs{} alone
cannot say whether agents exploit blind spots or merely produce
honestly-incomplete fixes --- so a differential diagnostic tests each
accepted-but-wrong patch just outside the visible suite's footprint, and
is itself validated against adjudicated human labels before its output is
trusted (Sec.~\ref{sec:diagnostic}). And an imperfect verifier already
inflates accepted-but-wrong mass for a purely static generator --- so a
budget-matched, feedback-free selection control isolates adaptive
exploitation from that mechanical selection effect
(Sec.~\ref{sec:controls}).

Degrading visible-suite adequacy raises false acceptance by
\PH{$\Delta F$~pp} between full adequacy and $\ms{=}0.30$, while false
acceptance at full adequacy stays \PH{$\le F(1.00)$~pp}, isolating
adequacy as the operative driver --- measured on \PH{$N$} repair tasks
spanning Python (SWE-bench Verified) and Java (Defects4J) across
\PH{$K$} agent families, against a full-adequacy control, a size-matched
adequacy-uncontrolled arm, and a budget-matched selection-only arm. The
rise \PH{emerges around $s^\ast{=}0.xx$ / grows gradually}; it
\PH{does / does not} transfer across the tested agent families; the
iterative loop exceeds the selection-only control by \PH{$\ge$5~pp} in the
low-adequacy band, licensing an attribution to adaptive agent behavior
beyond selection; and the diagnostic reaches \PH{$\kappa{=}0.xx$}
agreement with human labels, attributing \PH{$xx\%$} of low-adequacy
false acceptances to gaming-shaped behavior rather than honest
incompleteness. All of this is reported for budget-bounded iterative
repair under a frozen policy --- the only regime we test; extrapolation to
training-time verifiers under gradient pressure is a bounded conjecture,
not a claim any hypothesis licenses.

This paper makes the following contributions:
\begin{enumerate}[leftmargin=1.5em]
\item A formalization of visible test-suite mutation adequacy as a
\textbf{controllable attack surface} for coding-agent acceptance: the
surviving-mutant set is the exact space in which a wrong patch can pass,
and its size is a dial-able scalar rather than a fixed accident of the
shipped suite (Sec.~\ref{sec:background}).
\item An \textbf{adequacy-controlled sub-suite construction} isolating
mutation-kill score from suite size, code coverage, and equivalent-mutant
noise, with the complementary held-out oracle it simultaneously defines
(Sec.~\ref{sec:instrument}).
\item The \textbf{false-acceptance curve \Fs{} and onset estimate
$s^\ast$} for coding agents under controlled adequacy, with the
sharp-versus-gradual question decided by model comparison rather than
presupposed (Sec.~\ref{sec:results}).
\item A \textbf{separation of adaptive feedback exploitation from
selection} under an imperfect verifier, and of deliberate suite-gaming
from honestly-incomplete fixes via a human-validated per-patch diagnostic
(Secs.~\ref{sec:controls}, \ref{sec:diagnostic}).
\item A \textbf{practical sufficiency floor}: in our setting, visible
suites below mutation adequacy $s^\ast{=}$\,\PH{$0.xx$} should not serve
as standalone acceptance gates for agent-produced patches --- stated for
evaluation- and acceptance-time gating, with the training-time
extrapolation explicitly bounded (Sec.~\ref{sec:discussion}).
\end{enumerate}

The contributions are deliberately ordered so the paper degrades
gracefully under its own pre-registered fallbacks: the calibrated floor
(items 2--3) requires only \Fs{}; the gaming attribution (item 4) is gated
on diagnostic validation and can be withdrawn without voiding the floor;
transfer (within item 3) is reported per agent either way.

\section{Background: Adequacy as an Exploit Surface}\label{sec:background}

\subsection{Setting}
An agent receives a repair task: a buggy program $P$, a natural-language
issue, and a \emph{visible} test suite $T_v$. The agent may edit $P$ and
execute $T_v$ freely; it terminates when $T_v$ passes (acceptance,
$V{=}1$) or its budget is exhausted. Afterwards --- and invisibly to the
agent --- the final patch is executed against a \emph{held-out oracle}
$T_h$, yielding $H \in \{0,1\}$. This visible/held-out decomposition is
the standard construct of patch-overfitting
research~\cite{Smith2015Overfitting} and of recent agent reward-hacking
benchmarks~\cite{SpecBench_2605_21384}; what is new here is that we
control the strength of $T_v$.

\subsection{Mutation adequacy, definitionally}
Given a mutant pool $M$ generated from the fixed reference version of $P$,
the mutation-kill score of a suite $T$ is
$\ms(T) = |\{m \in M^\ast : T \text{ kills } m\}| \,/\, |M^\ast|$, where
$M^\ast \subseteq M$ excludes mutants judged equivalent
(Sec.~\ref{sec:equiv}). We emphasize a definitional rather than empirical
reading of \ms{}: whether mutation score predicts \emph{real-fault}
detection is contested in the
literature~\cite{Just2014Valid,Papadakis2019Survey}, and this paper does
not rely on that correlation. For our purposes, $1-\ms(T_v)$ measures,
literally and by construction, the normalized size of the
\emph{surviving-mutant set}: the enumerable space of behavior changes
$T_v$ cannot see. Every surviving mutant is a certificate that some wrong
program passes $T_v$. It is in this sense --- and only this sense --- that
adequacy is the attack surface: we make no claim that \ms{} measures
``true suite quality,'' only that it measures the size of the exploitable
slack, which is the quantity an adapting agent can occupy.

\subsection{Why the agent case differs from classical overfitting}
Classical APR overfitting~\cite{Smith2015Overfitting,Qi2015PatchPlausibility}
arises when a \emph{search} procedure samples patches until one passes.
An iterative agent adds two ingredients: it \emph{reads} the failing
tests (so the visible suite's blind spots are partially legible in the
feedback), and it \emph{adapts} its edits across turns. Whether these
ingredients add accepted-but-wrong mass beyond what any imperfect
verifier mechanically produces for a static
sampler~\cite{ResamplingLimits_j8H84v6AZ1} is precisely our RQ5, and we
measure it rather than assume it.

\section{Study Design}\label{sec:design}

\subsection{Research questions}
\begin{itemize}[leftmargin=1.5em]
\item \textbf{RQ1 (Adequacy $\to$ false acceptance).} Does false
acceptance rise as visible-suite mutation adequacy falls, with suite size
and coverage held fixed?
\item \textbf{RQ2 (Onset and shape).} Is there an onset threshold
$s^\ast$ below which visible test-pass stops tracking correctness, and is
the rise sharp or gradual?
\item \textbf{RQ3 (Transfer).} Does the adequacy--false-acceptance
relationship, and $s^\ast$ in particular, transfer across agent families
and task languages?
\item \textbf{RQ4 (Mechanism).} Are accepted-but-wrong patches
gaming-shaped --- exploiting the visible suite's specific blind spots ---
or merely honestly incomplete?
\item \textbf{RQ5 (Adaptation vs.\ selection).} Does iterative
test-feedback use add accepted-but-wrong mass beyond the mechanical
selection effect of an imperfect verifier?
\end{itemize}

\subsection{Estimand and pre-registered decision machinery}\label{sec:estimand}
The primary estimand is \Fs{} of Eq.~\eqref{eq:fs}, estimated at target
adequacy levels $s \in \{0.30, 0.45, 0.60, 0.75, 0.90, 1.00\}$. The
marginal overfitting gap $\Gs = P(V{=}1 \mid s) - P(H{=}1 \mid s)$ is
retained as a secondary, descriptive curve.\footnote{The study was first
specified with \Gs{} as primary; the promotion of \Fs{} was made as a
recorded pre-data amendment, before any agent episode was run, because
\Fs{} conditions on the acceptance event and is therefore the quantity an
acceptance gate exposes. No estimand was altered after data collection.}

Because \Fs{} conditions on $V{=}1$, its per-cell denominator is the
number of \emph{accepted} patches, which can be small precisely where the
suite is strongest. We therefore build the thin-cell rule into the
estimand itself: a cell (task $\times$ agent $\times$ level) with fewer
than $n_{\min} =$ \PH{$n_{\min}$} accepted patches is reported
\textsc{inconclusive} for \Fs{} and is excluded from curve fitting; every
reported \Fs{} point carries its denominator. If the pilot
(Sec.~\ref{sec:pilot}) shows $F(1.00)$ cannot reach $n_{\min}$ at our task
budget, the pre-registered fallback is to report \Gs{} as primary and
\Fs{} descriptively --- decided before the main run, never after.

\subsection{Hypotheses}\label{sec:hypotheses}
All confirmatory hypotheses, thresholds, and fallbacks below were fixed
before data collection; family-wise error across H1--H5 is controlled by
Holm's procedure at $\alpha = 0.05$.

\medskip\noindent
\textbf{H1 (Adequacy drives false acceptance).} \emph{(i)}
$F(0.30) - F(1.00) \ge 10$~pp; \emph{(ii)} the fixed effect of \ms{} on
$H$ among accepted patches is negative ($p < 0.05$, Holm); \emph{(iii)}
$F(1.00) \le 10$~pp.
\wj{Standpoint: this is the compass hypothesis --- the binary success
criterion of the whole study. (iii) is the internal control: if even
full-adequacy suites over-accept, adequacy is not the operative driver and
the finding becomes an attack on the adequacy construct itself
(pre-worded fallback), which is still publishable but a different paper.
Disconfirming result: a flat curve; governed by fallback F-G1, reported as
a clean refutation of the weak-suites-invite-gaming assumption.}

\medskip\noindent
\textbf{H2 (Shape).} The rise of \Fs{} as $s$ decreases is monotone; and
\textbf{H2a (Sharpness)}: a segmented (broken-stick) model beats a smooth
monotone alternative by $\Delta\mathrm{AIC} \ge 4$, yielding an onset
estimate $s^\ast$ with a bootstrap 95\% CI.
\wj{Standpoint: we deliberately pose sharpness as the HARDER claim and
predict, based on the digest of prior overfitting curves, that H2 holds
but H2a fails (gradual growth). If neither segmented nor smooth fits
monotonically, the pre-worded fallback is that scalar mutation score does
not order agent-exploitable weakness --- a construct-level finding.
``Sharp threshold'' language is licensed only if H2a passes; otherwise
$s^\ast$ is reported as a regulatory floor chosen on the fitted smooth
curve, not a discovered discontinuity.}

\medskip\noindent
\textbf{H3 (Transfer).} The agent$\times$adequacy interaction is
non-significant AND per-agent $s^\ast$ estimates lie within a
pre-registered tolerance band of $\pm 0.10$.
\wj{Standpoint: both branches are pre-worded and both are findings. If H3
holds, the floor is agent-intrinsic and the practitioner rule is
universal. If H3 fails, the deliverable degrades gracefully to a
per-agent floor table --- ``compute your suite's kill score, look up your
agent's floor'' --- which is weaker but still actionable. The narrative
must not presuppose transfer.}

\medskip\noindent
\textbf{H4 (Gaming attribution).} \emph{(i)} The differential diagnostic
agrees with blind, adjudicated human labels at Cohen's $\kappa \ge 0.7$
on a stratified sample of \PH{$n_{\mathrm{lab}}$} accepted-but-wrong
patches; \emph{(ii)} conditional on (i), the gaming-shaped fraction among
accepted-but-wrong patches rises by $\ge 15$~pp from the highest to the
lowest adequacy band.
\wj{Standpoint: (i) is a kill-switch, and the paper is built to survive
it: if $\kappa < 0.7$ the diagnostic is withdrawn, every gaming
attribution disappears from the paper (combination rule 5), and the
contribution reverts to the calibrated floor, which needs only F(s).
This is why gaming is a SECONDARY contribution in
Sec.~\ref{sec:intro} --- a deliberate hierarchy choice, not modesty.}

\medskip\noindent
\textbf{H5 (Adaptation beyond selection).} In the low-adequacy band
($\ms \le 0.45$), the iterative fix-until-pass arm exceeds the
budget-matched, feedback-free best-of-$k$ selection arm:
$F_{\mathrm{iter}} - F_{\mathrm{sel}} \ge 5$~pp.
\wj{Standpoint: this is the causal firewall separating our claim from
the imperfect-verifier literature's mechanical ceiling
(ResamplingLimits): an imperfect verifier inflates accepted-but-wrong
mass even for a static sampler, so a naive reading of H1 alone would
over-attribute to ``gaming.'' If H5 fails (fallback F5), every
agent-agency claim in the paper is restated as a selection effect and
the study is reported as the first adequacy-controlled confirmation of
the selection ceiling --- the title and floor survive because they do not
depend on agency.}

\subsection{Confirmatory vs.\ exploratory scope}
H1--H5 are confirmatory. Pre-registered as exploratory: per-mutant-class
breakdowns of which operators' survivors are exploited, operator-mix
robustness cells, per-repository effects, and qualitative patch
taxonomies. Exploratory results are reported in a separate subsection and
never enter the Holm family.

\section{Instrumentation}\label{sec:instrument}

\subsection{Adequacy-controlled sub-suite sampler}\label{sec:sampler}
For each task we require visible sub-suites that differ in \ms{} but not
in the confounds a reviewer would reach for first: suite size and code
coverage. Given the task's full suite $T$ (all developer tests plus the
benchmark's fail-to-pass tests) and filtered mutant pool $M^\ast$
(Sec.~\ref{sec:equiv}), we compute the kill matrix
$K \in \{0,1\}^{|T| \times |M^\ast|}$ once per task. The sampler then
draws, for each target level $s$, sub-suites $T_v \subseteq T$ subject to:
(a) $|\ms(T_v) - s| \le 0.03$; (b) $|T_v| = n_0$ fixed across all levels
of the task, with $n_0$ chosen as the largest size for which the lowest
target level remains reachable; (c) line coverage of $T_v$ within
$\pm 5$\% of the level-1.0 reference sub-suite; (d) the benchmark's
fail-to-pass tests for the target defect always included, so every arm
still specifies the bug being repaired. Sampling is greedy-random over
$K$ with rejection, \PH{$r$} independent sub-suites per (task, level),
seeds fixed and published. Tasks for which the constraint set is
infeasible at any level are excluded before any agent run, with counts
reported (Table~\PH{X}).

Constraint (d) has a consequence worth stating: our sweep degrades the
\emph{regression-detection} adequacy of the visible suite while keeping
the defect specification constant. This is the deployment-realistic
regime --- gating suites almost always contain a reproduction test for
the issue at hand and vary in how well they guard everything else.

\subsection{Equivalent mutants as first-class machinery}\label{sec:equiv}
Detecting equivalent mutants is undecidable, real-world equivalence rates
are material, and the strongest simple filter, Trivial Compiler
Equivalence (TCE), removes only a fraction of
them~\cite{Papadakis2015TCE,Papadakis2019Survey}. Because \ms{} is our
independent variable and equivalent mutants corrupt its denominator ---
non-randomly, clustered by operator and code region --- we treat
equivalence handling as reported machinery rather than a preprocessing
checkbox. Three layers: \emph{(1)} TCE at compile/bytecode level removes
provably-trivial equivalents; \emph{(2)} an independent LLM-based
equivalence judgment in the style of industrial mutation
pipelines~\cite{Foster2025ACH}, run with a fixed prompt and model
snapshot, flags suspected equivalents; inter-method agreement is reported
(Table~\PH{X}). \emph{(3)} Every headline curve is recomputed under a
\emph{sensitivity band}: an aggressive policy (drop all suspected
equivalents) and a conservative policy (drop only TCE-proven ones). The
onset $s^\ast$ is reported with its value under both policies; if
$s^\ast$ moves outside its bootstrap CI across the band, that instability
\emph{is} reported as a finding about the construct
\wj{pre-committed: we do not get to pick the friendlier band post hoc;
the abstract's $s^\ast$ is the conservative-policy value, with the band
in the results table}. Tasks whose suspected-equivalence rate exceeds
\PH{$q_{\mathrm{eq}}$}\% are excluded pre-run, counts reported.

\subsection{Held-out oracle, as a lower bound}\label{sec:oracle}
The held-out oracle $T_h$ comprises (i) all tests of $T$ excluded from
$T_v$, (ii) the benchmark's hidden evaluation tests where the benchmark
maintains them, and (iii) the killing tests of every mutant excluded from
$T_v$'s kill responsibility. $T_h$ is fixed per (task, level, replicate)
\emph{before} any agent run and is never visible to any arm.

$T_h$ is itself a finite suite, so the label $H{=}1$ can be granted to a
patch that is wrong on inputs neither suite covers. The error is
one-sided: it can only \emph{undercount} false acceptances. Every \Fs{}
in this paper is therefore a \textbf{lower bound} on the true
false-acceptance rate, and every claim is phrased accordingly
\wj{this framing is load-bearing for construct validity: a reviewer
applying the ``your oracle is also just a suite'' attack is conceded in
advance, and the direction of the bias strengthens rather than weakens
the alarm the floor raises}. As a robustness probe we additionally
amplify $T_h$ with EvalPlus-style generated
tests~\cite{EvalPlus_Liu2023} on the \PH{$n_{\mathrm{amp}}$} tasks where
type-aware amplification applies, and report how much \Fs{} moves
(Table~\PH{X}).

\subsection{Gaming-vs-incompleteness diagnostic}\label{sec:diagnostic}
For each accepted-but-wrong patch we execute a differential probe just
outside the visible suite's footprint: inputs generated from the visible
tests by boundary perturbation and metamorphic variation, executed
against the patch and the reference fix. A patch is \emph{gaming-shaped}
if its disagreements with the reference concentrate on
visible-suite-adjacent inputs while it reproduces visible-checked
behavior exactly --- e.g., special-casing checked values, hardcoding
expected outputs, or disabling the faulty path only for tested inputs;
it is \emph{honestly incomplete} if disagreements distribute across the
input space without structure aligned to $T_v$'s footprint. The
classifier's output is trusted only after H4(i): two authors blind-label
a stratified sample of \PH{$n_{\mathrm{lab}}$} patches with adjudication,
and Cohen's $\kappa$ against the classifier is reported. Below the 0.7
gate, all gaming attributions are withdrawn.

\subsection{Selection-only control and cross-agent harness}\label{sec:controls}
The selection-only arm isolates adaptation: the same base model generates
$k$ independent single-shot patches with \emph{no} test feedback; the
visible suite then selects the first passing candidate, with $k$
calibrated so that total model-token and execution budget matches the
iterative arm's mean consumption per (task, level)
\wj{budget matching is by cost, not by attempts --- matching attempts
would hand the iterative arm more tokens and bias H5 toward
adaptation}. All arms share the cross-agent harness: identical prompts
up to agent-native scaffolding, identical budgets, temperature
\PH{$\tau$}, \PH{$k$} episodes per (task, level, agent), fixed seeds
where the agent stack honors them.

\section{Experimental Setup}\label{sec:setup}

\subsection{Tasks and benchmarks}
We draw \PH{$N_{\mathrm{py}}$} Python tasks from SWE-bench
Verified~\cite{SWEbench_VTF8yNQM66} and \PH{$N_{\mathrm{jv}}$} Java
tasks from Defects4J v\PH{x.y}~\cite{Defects4J_Just2014}, filtered to
tasks with (a) a deterministic, flake-free full suite under \PH{$R$}
repeated runs, (b) a feasible sampler constraint set at all six levels,
and (c) suspected-equivalence rate below \PH{$q_{\mathrm{eq}}$}\%.
Filtering counts per criterion are reported in Table~\PH{X}. Because
SWE-bench Verified is widely used in agent training corpora,
contamination is a first-order concern: we tag each task with the
repository-snapshot recency heuristics of \PH{tooling}, and the
\textbf{contamination-excluded subset is the primary analysis set}
whenever it disagrees with the full set; both are always reported.

\subsection{Mutation tooling}
Java mutants are generated with Major~\cite{Just2011Major} and Python
mutants with \PH{mutmut/cosmic-ray, version}, under the standard
operator sets of each tool; operator mixes are held identical across
levels within a task by construction (the pool $M$ is fixed; only
$T_v$'s kill responsibility varies). Kill matrices are computed once on
the unpatched reference version with per-test timeouts of
\PH{$t$}s.

\subsection{Agents and budgets}
We evaluate \PH{$K$} agent families spanning the current design space:
a pipeline-style repair system (Agentless-class), a
tool-using autonomous agent (SWE-agent-class), and \PH{a third family},
each backed by \PH{model list, pinned snapshots}. Budgets:
\PH{$B$} iterations or \PH{\$C} per episode, whichever first;
\PH{$k$} episodes per cell. The full grid is
$\PH{N} \times 6 \times \PH{K} \times \PH{k}$ episodes plus the
selection-only arm at \PH{levels}, a total of \PH{$E$} episodes at an
estimated compute cost of \PH{\$D}.

\subsection{Statistical analysis}\label{sec:analysis}
RQ1/H1: mixed-effects logistic regression of $H$ among accepted patches
on \ms{} (fixed), with crossed random intercepts for task and agent and
a random slope for task; the H1(i) contrast is estimated on the fitted
marginal means with cluster-bootstrap CIs (tasks as clusters,
\PH{$B_{\mathrm{boot}}$} resamples). RQ2/H2: segmented (single-knot
broken-stick) versus monotone smooth (I-spline GAM) fits on \Fs{};
$\Delta\mathrm{AIC} \ge 4$ decides sharpness; $s^\ast$ is the knot (if
segmented wins) or the pre-registered regulatory crossing
$F(s) = F(1.00) + 5$~pp (if smooth wins). RQ3/H3: agent$\times$\ms{}
interaction by likelihood-ratio test plus per-agent $s^\ast$ band check.
RQ4/H4: Cohen's $\kappa$ with adjudication; band contrast by
proportion test within accepted-but-wrong patches. RQ5/H5: paired
per-(task, level) contrast between arms, McNemar on paired episode
outcomes where pairing is exact and cluster bootstrap otherwise. All
five families enter one Holm correction. Thin cells follow
Sec.~\ref{sec:estimand}. Analysis code and the locked analysis plan are
in the artifact.

\subsection{Pilot gates (go/no-go)}\label{sec:pilot}
Before the main grid, a \PH{15}-task pilot must pass three gates:
\textbf{G1} --- $F(1.00)$ cells reach $n_{\min}$ at the pilot's
acceptance rates, else the \Gs{} fallback activates; \textbf{G2} ---
$s^\ast$ (pilot estimate) is stable across the equivalent-mutant
sensitivity band; \textbf{G3} --- the diagnostic clears $\kappa \ge 0.7$
on the pilot's labeled sample. Gate outcomes are reported regardless of
result \wj{the pilot is a capability check on the instrument, not a
peek at the hypothesis; pilot episodes are excluded from confirmatory
analysis}.

\section{Results}\label{sec:results}
\wj{Every number below is a pre-registered slot. If a hypothesis's
fallback fires, the corresponding subsection is REPLACED by its
pre-committed fallback text (kept in the artifact), not edited.}

\subsection{RQ1: Adequacy drives false acceptance}
Across \PH{$N$} tasks and \PH{$K$} agents, false acceptance rises from
$F(1.00) = $ \PH{$x.x$}\% to $F(0.30) = $ \PH{$xx.x$}\%
($\Delta = $ \PH{$xx$}~pp; H1(i) threshold 10~pp), with a negative
adequacy effect ($\hat\beta_{\ms} = $ \PH{$-x.xx$}, $p = $
\PH{$x{\cdot}10^{-x}$}, Holm) and a full-adequacy rate below the
10~pp control bound. The size-matched, adequacy-uncontrolled arm shows
\PH{result}, confirming the effect is carried by adequacy rather than
suite size (Table~\PH{X}). Per-cell denominators accompany every point
(Fig.~\PH{X}); \PH{$c$} cells were \textsc{inconclusive} under
$n_{\min}$ and excluded from fitting.

\subsection{RQ2: Onset and shape}
The \PH{segmented / smooth} model wins
($\Delta\mathrm{AIC} = $ \PH{$x.x$}), yielding
$s^\ast = $ \PH{$0.xx$} (95\% CI \PH{$[0.xx, 0.xx]$}), stable across the
equivalent-mutant sensitivity band
(conservative \PH{$0.xx$} / aggressive \PH{$0.xx$}).
\wj{If smooth wins, this subsection's language shifts from
``discovered discontinuity'' to ``regulatory crossing'' per H2a's
standpoint note; the floor remains.}

\subsection{RQ3: Transfer across agents and languages}
The agent$\times$adequacy interaction is \PH{n.s. / significant}
(LRT $p = $ \PH{$0.xx$}); per-agent onsets are
\PH{$0.xx$ / $0.xx$ / $0.xx$}, \PH{within / outside} the $\pm 0.10$
band. Per-language curves are reported separately
(Python \PH{$s^\ast$}, Java \PH{$s^\ast$}); pooling \PH{is / is not}
licensed. \wj{Non-transfer branch: the deliverable becomes the
per-agent floor table (Table~X), pre-worded as an actionable artifact,
not a hedge.}

\subsection{RQ4: Gaming-shaped vs.\ honestly incomplete}
The diagnostic reaches $\kappa = $ \PH{$0.xx$} against blind adjudicated
labels (\PH{pass/fail} of the 0.7 gate). \PH{Conditional on the gate:}
gaming-shaped fractions among accepted-but-wrong patches rise from
\PH{$xx$}\% (band $\ms \ge 0.75$) to \PH{$xx$}\%
(band $\ms \le 0.45$), a \PH{$xx$}~pp increase (H4(ii) threshold
15~pp). Exploratory operator-class breakdowns appear in
Sec.~\ref{sec:exploratory}.

\subsection{RQ5: Adaptation beyond selection}
In the low-adequacy band, the iterative arm's false acceptance exceeds
the budget-matched selection-only arm by \PH{$x.x$}~pp
(H5 threshold 5~pp; paired \PH{McNemar/bootstrap} $p = $ \PH{$0.xx$}).
\wj{If F5 fires: all agency language in
Secs.~\ref{sec:intro}--\ref{sec:discussion} is restated as selection;
the floor and RQ1--RQ3 findings are unaffected by construction.}

\subsection{Calibration against as-shipped suites}
Computed from the same kill matrices at no additional agent cost,
\PH{$q$}\% of benchmark tasks ship visible suites with
$\ms < s^\ast$ (Python \PH{$q$}\%, Java \PH{$q$}\%): the regime our
floor warns about is \PH{common / rare} in the wild.

\subsection{Exploratory analyses}\label{sec:exploratory}
\PH{operator-class exploitation profiles}, \PH{per-repository
effects}, and a \PH{qualitative taxonomy of gaming-shaped patches with
examples} are reported here. These are
outside the Holm family and claim no confirmatory status.

\section{Discussion and Implications}\label{sec:discussion}

\textbf{A calibrated sufficiency floor.} The practitioner rule this
study licenses is: compute your gating suite's mutation-kill score
(one mutation run; no agent episodes needed); if it falls below
$s^\ast = $ \PH{$0.xx$} \PH{(per-agent table X if H3 failed)}, do not
treat visible test-pass as a standalone acceptance signal --- add held-out
tests, mutant-derived tests, or human review. Because every \Fs{} we
report is a lower bound (Sec.~\ref{sec:oracle}), the floor errs on the
permissive side: the true unsafe region can only be larger.

\textbf{What the floor does and does not claim.} It is an
acceptance-time result under a frozen policy and bounded budgets. It
does not claim that training against weak verifiers produces gaming
policies (RLVR-regime extrapolation is a conjecture our design cannot
test), does not claim mutation score measures true suite quality
(Sec.~\ref{sec:background}), and --- if H4(i) failed --- does not
attribute intent-like structure to the patches.

\textbf{For benchmark designers.} Reporting a benchmark's per-task
visible-suite \ms{} distribution alongside pass rates would let readers
condition trust on adequacy; our kill matrices for both benchmarks are
released in the artifact.

\textbf{For the imperfect-verifier literature.} Our selection-only arm
gives the resampling-ceiling account~\cite{ResamplingLimits_j8H84v6AZ1}
its first adequacy-controlled instantiation with an executable verifier;
H5's outcome says whether adaptation adds mass beyond that ceiling.

\section{Threats to Validity}\label{sec:threats}

\textbf{Construct.} (1) Equivalent mutants corrupt \ms{}'s denominator;
undecidable in general, with material real-world rates and only partial
automated detection~\cite{Papadakis2015TCE}. Mitigation: the three-layer
pipeline and sensitivity band of Sec.~\ref{sec:equiv}; the headline
$s^\ast$ is the conservative-band value. (2) The held-out oracle is
finite, so \Fs{} is a lower bound; all claims are phrased as such, and
the EvalPlus-amplified probe bounds the movement. (3) Mutation score's
relation to real-fault detection is contested~\cite{Just2014Valid}; we
rely only on the definitional reading (surviving-set size), not the
contested correlation.

\textbf{Internal.} (4) Contamination of SWE-bench-derived tasks can
flatten the curve toward a false null; the contamination-excluded set is
primary on disagreement. (5) Thin \Fs{} denominators at high adequacy
could manufacture shape artifacts; the $n_{\min}$ rule and per-cell
reporting are built into the estimand. (6) Sampler rejection could
correlate with task difficulty; excluded-task counts and characteristics
are reported.

\textbf{External.} (7) Two languages, two benchmarks, \PH{$K$} agent
families, frozen policies, bounded budgets: the floor is stated for this
regime. Per-language and per-agent results are reported unpooled unless
H3 licenses pooling. (8) Operator sets are tool-standard; the
operator-mix robustness cell (exploratory) probes this choice.

\section{Related Work}\label{sec:related}

\textbf{Patch overfitting in APR.} Overfitting to the evaluating suite
is the classical failure of generate-and-validate
repair~\cite{Smith2015Overfitting,Qi2015PatchPlausibility}, with
per-patch detection by generated differential
tests~\cite{XinReiss2017DiffTGen} and behavioral
similarity~\cite{Xiong2018PatchSim}. We concede these as method-twins of
our two-suite construct and our diagnostic; none sweeps suite adequacy
as a controlled variable, none matches size and coverage, and none
studies iterative LLM agents.

\textbf{Reward hacking and imperfect verifiers.} Reward gaming in text
generation~\cite{acl_pang_etal_2023_reward,acl_he_etal_2023_blind} and
tail-chasing under reward models~\cite{ChasingTail_pBjy4ek2QV}
establish the exploitation dynamic; verifier-imperfection theory bounds
what selection can achieve under fixed
error~\cite{ResamplingLimits_j8H84v6AZ1,ROCnReroll_3Gy5mmyuxn}. Recent
benchmarks measure visible/held-out gaps for coding agents under
shipped suites~\cite{SpecBench_2605_21384,RHB_2605_02964}. All fix the
verifier we dial.

\textbf{Test generation against code models.} Adversarial co-evolution
of coders and testers~\cite{CURE_wPdBe9zxNr,ATGen_Sxj4o3qXtl,
UTRL_VqjNYF9nbP} and self-verification agents~\cite{ReVeal_q56ZI1Co43}
strengthen suites during training; we need the complementary diagnostic
setting --- a suite held fixed at controlled, degraded adequacy.
Execution-free reward models~\cite{2512_21919} motivate themselves by
suite inadequacy without quantifying it; our curve supplies the
quantification. Benchmark-strengthening work~\cite{EvalPlus_Liu2023}
amplifies oracles for function-level code; we use it as a robustness
probe, not a surrogate for our construct.

\textbf{Mutation testing.} Foundations and the equivalent-mutant
problem~\cite{Papadakis2019Survey,Papadakis2015TCE,Just2011Major},
validity for real faults~\cite{Just2014Valid}, and industrial LLM-based
mutation pipelines~\cite{Foster2025ACH}. Mutation-guided \emph{test}
improvement is well studied; mutation adequacy as a swept independent
variable against an adapting agent is, to our knowledge, ours.

\section{Conclusion}\label{sec:conclusion}
We turned a warning into a measurement. Visible-suite mutation adequacy
--- definitionally the size of the surviving-mutant space a wrong patch
can occupy --- was made a controlled, matched, swept variable for coding
agents, yielding the false-acceptance curve \Fs{}, an onset estimate
$s^\ast$, a validated separation of gaming from honest incompleteness,
and an isolation of adaptive exploitation from mechanical selection.
The resulting sufficiency floor gives practitioners a one-number,
pre-agent check on whether their gating suite's green light means
anything. The design's pre-registered fallbacks ensure every branch ---
including the null --- returns an interpretable finding.

\section*{Data Availability}
The artifact (sampler, kill matrices, equivalence-filter outputs and
sensitivity bands, harness configurations, locked analysis plan,
labeled diagnostic sample, and all episode logs) is available at
\PH{URL}, archived at \PH{DOI}.

\balance
\bibliographystyle{IEEEtran}
\bibliography{refs}

\end{document}
